% from main.tex
%%%%%%%%%%  scan idx 619  |  folio 606  %%%%%%%%%%
En
effet on peut évidemment supposer \(S\) et \(T\) affines. Prenons un recouvrement affine \(\{V_i\}\) de \(X_T\). Soient \(g\in\Gamma(X_T;\mathcal O_{X'_T}^{*})\)~; \(g_i\) un relèvement de \(g|V_i\) dans \(\Gamma(V_i;\mathcal O_{X_T}^{*})\). Alors \(\delta_T^*(g)\) est égal à la classe du \(1\)-cocycle \(\delta_{ij}^*=(g_ig_j^{-1}-1)\)~; \(\delta_T(g)\), à celle de \(\delta_{ij}=(g_i-g_j)=g_j\delta_{ij}^*=g\,\delta_{ij}^*\), donc on a
\[
  \delta_T(g)=g\,\delta_T^*(g)
\]
en tenant compte de l'opération de \(\Gamma(X_T,\mathcal O_{X'_T})\) sur \(I_T\), donc sur \(H^1(X_T,I_T)\).

D'après ii), on a \(\Gamma(X_T;\mathcal O_{X'_T})=\Gamma(X;\mathcal O_{X'})\otimes_S\Gamma(T,\mathcal O_T)\) et donc \(g=\sum g_\beta\otimes t_\beta\) pour
\(g_\beta\in\Gamma(X,\mathcal O_{X'})\) et
\(t_\beta\in\Gamma(T,\mathcal O_T)\). Or par (v) \((f_*\mathcal O_X)_{\red}\longrightarrow(f_*\mathcal O_{X'})_{\red}\) est surjectif~; par suite \(g_\beta=v(h_\beta)+n_\beta\) pour
\(h_\beta\in\Gamma(X;\mathcal O_X)\), \(v:\Gamma(X,\mathcal O_X)\to
\Gamma(X,\mathcal O_{X'})\) et
\(n_\beta\in\operatorname{nilrad}(\Gamma(X,\mathcal O_{X'}))\). D'après i),
\((n_\beta\otimes t_\beta)I_T=0\). Par suite \(\delta_T(g)=h\,\delta_T^*(g)\) où \(h=\sum v(h_\beta)\otimes t_\beta\in\Gamma(X_T,\mathcal O_{X_T})\) d'où \(\delta_T^*(g)=h^{-1}\delta_T(g)=\delta_T(h^{-1}g)\), donc \(\im\delta_T^*\subset\im\delta_T\).

Finalement soit \(g=\sum g_\beta\otimes t_\alpha\in
\Gamma(X_T,\mathcal O_{X'_T})\). Soit \(g_\beta=v(h)+n_\beta\) où
\(n_\beta\in\operatorname{nilrad}(\Gamma(X,\mathcal O_X))\); on a alors \(\delta_T(g)=\delta_T\bigl(\sum n_\beta\otimes t_\beta\bigr)=\delta_T\bigl(1+\sum n_\beta\otimes t_\beta\bigr)=\delta_T^*\bigl(1+\sum n_\beta\otimes t_\beta\bigr)\), d'où \(\im\delta_T\subset\im\delta_T^*\).

En vertu de 3.2, iii) et iv), on a une suite exacte fonctorielle en \(T\)~:
\[
  0\longrightarrow \mathbb V(E^\vee)(T)\longrightarrow
  R^1(f_t)_*(\mathcal O_{X_T}^{*})
  \longrightarrow
  R^1(f_T)_*(\mathcal O_{X'_T})
  \longrightarrow
  \mathbb V(Q)(T),
\]
qui, par passage aux faisceaux f.p.p.f. associés, donne une suite exacte
\[
  0\longrightarrow \mathbb V(E^\vee)\longrightarrow\Pic_{X/S}
  \longrightarrow\Pic_{X'/S}\xrightarrow{u}\mathbb V(Q).
\]
Si \(P=\Ker(u)\), le monomorphisme \(P\hookrightarrow\Pic_{X'/S}\) s'obtient par changement de base à partir de la section unité \(S\hookrightarrow\mathbb V(Q)\)~; comme \(\mathbb V(Q)\) est représentable et séparé, il résulte que \(P\hookrightarrow\Pic_{X'/S}\) est représentable par une
